\bta{2003}


\begin{enumerate}
	%\renewcommand{\labelenumi}{\arabic{enumi}.}
	% A(\Alph) a(\alph) I(\Roman) i(\roman) 1(\arabic)
	%设定全局标号series=example	%引用全局变量resume=example
	%[topsep=-0.3em,parsep=-0.3em,itemsep=-0.3em,partopsep=-0.3em]
	%可使用leftmargin调整列表环境左边的空白长度 [leftmargin=0em]
	\item
\begin{enumerate}
	%\renewcommand{\labelenumi}{\arabic{enumi}.}
	% A(\Alph) a(\alph) I(\Roman) i(\roman) 1(\arabic)
	%设定全局标号series=example	%引用全局变量resume=example
	%[topsep=-0.3em,parsep=-0.3em,itemsep=-0.3em,partopsep=-0.3em]
	%可使用leftmargin调整列表环境左边的空白长度 [leftmargin=0em]
	\item
	如果厄密算符 $\hat{A}$ 对任何矢量 $|u\rangle$, 有 $\langle u|\hat{A}| u\rangle \geq 0$, 则称 $\hat{A}$ 为正定算符, 求证：算符 $\hat{A}=|a\rangle\langle a|$ 是厄密正定算符。
	\banswer{
		由 $\hat A=|a\rangle\langle a|$ 可得
		\[
			\hat A^\dagger=(|a\rangle\langle a|)^\dagger
			=|a\rangle\langle a|=\hat A,
		\]
		故 $\hat A$ 是厄米算符。对任意 $|u\rangle$，
		\[
			\langle u|\hat A|u\rangle
			=\langle u|a\rangle\langle a|u\rangle
			=|\langle a|u\rangle|^2\geq 0.
		\]
		因此，按题目给出的定义，$|a\rangle\langle a|$ 是正定算符（更严格地说是半正定算符；与 $|a\rangle$ 正交的非零态使平均值为零）。
	}
	
	\item 
	如果 $\hat{A}$ 是任 一线性算符, 求证 $A^{+} A$ 是正定的厄密算符, 它的迹等于 $\hat{A}$ 在任意表象中的矩阵元的模平方之和。试推导, 当且仅当 $\hat{A}=0$ 时, $\operatorname{Tr}\left(A^{+} A\right)=0$ 才成立。
	
	\banswer{
		首先
		\[
			(\hat A^\dagger\hat A)^\dagger
			=\hat A^\dagger\hat A,
		\]
		所以 $\hat A^\dagger\hat A$ 是厄米算符；并且对任意 $|u\rangle$，
		\[
			\langle u|\hat A^\dagger\hat A|u\rangle
			=\langle \hat A u|\hat A u\rangle
			=\|\hat A|u\rangle\|^2\geq0,
		\]
		故其为正定（半正定）算符。

		在任意一组完备正交基 $\{|i\rangle\}$ 中插入恒等算符，得
		\begin{align*}
			\operatorname{Tr}(\hat A^\dagger\hat A)
			&=\sum_i\langle i|\hat A^\dagger\hat A|i\rangle
			=\sum_{i,j}\langle i|\hat A^\dagger|j\rangle
			\langle j|\hat A|i\rangle\\
			&=\sum_{i,j}|\langle j|\hat A|i\rangle|^2.
		\end{align*}
		右端是非负数之和。因此 $\operatorname{Tr}(\hat A^\dagger\hat A)=0$ 当且仅当所有矩阵元 $\langle j|\hat A|i\rangle$ 均为零，也即当且仅当 $\hat A=0$。
	}
	
	
	\item 
	求证：如果 $[[\hat{A}, \hat{B}], \hat{A}]=[[\hat{A}, \hat{B}], \hat{B}]=0$, 则 $e^{\hat{A}} e^{\hat{B}}=e^{\hat{A}+\hat{B}} e^{\frac{1}{2}[\hat{A}, \hat{B}]}$。
	\banswer{
		记 $\hat C=[\hat A,\hat B]$。题设给出 $[\hat C,\hat A]=[\hat C,\hat B]=0$。令
		\[
			F(t)=e^{t\hat A}e^{t\hat B}.
		\]
		由于 $e^{t\hat A}\hat B e^{-t\hat A}=\hat B+t[\hat A,\hat B]=\hat B+t\hat C$，有
		\begin{align*}
			\frac{dF}{dt}
			&=\hat A F+e^{t\hat A}\hat B e^{t\hat B}
			=(\hat A+\hat B+t\hat C)F.
		\end{align*}
		$\hat A+\hat B$ 与 $\hat C$ 对易，且 $F(0)=1$，故
		\[
			F(t)=\exp\!\left[t(\hat A+\hat B)+\frac{t^2}{2}\hat C\right].
		\]
		取 $t=1$，即得
		\[
			e^{\hat A}e^{\hat B}
			=e^{\hat A+\hat B+\frac12[\hat A,\hat B]}
			=e^{\hat A+\hat B}e^{\frac12[\hat A,\hat B]}.
		\]
	}
	
	\item 
	求证：任一可观测的平均值对时间的导数由下式给出：
	\[
	i \hbar \frac{d\langle\hat{A}\rangle}{d t}=\langle[\hat{A}, \hat{H}]\rangle+i \hbar\left\langle\frac{\partial \hat{A}}{\partial t}\right\rangle
	\]
	\banswer{
		设 $|\psi(t)\rangle$ 满足薛定谔方程
		$i\hbar|\dot\psi\rangle=\hat H|\psi\rangle$，则
		$\langle\dot\psi|=(i/\hbar)\langle\psi|\hat H$。于是
		\begin{align*}
			\frac{d\langle\hat A\rangle}{dt}
			&=\langle\dot\psi|\hat A|\psi\rangle
			+\left\langle\psi\left|\frac{\partial\hat A}{\partial t}\right|\psi\right\rangle
			+\langle\psi|\hat A|\dot\psi\rangle\\
			&=\frac{i}{\hbar}\langle\hat H\hat A\rangle
			-\frac{i}{\hbar}\langle\hat A\hat H\rangle
			+\left\langle\frac{\partial\hat A}{\partial t}\right\rangle\\
			&=\frac{1}{i\hbar}\langle[\hat A,\hat H]\rangle
			+\left\langle\frac{\partial\hat A}{\partial t}\right\rangle.
		\end{align*}
		两边乘以 $i\hbar$，便得到题中公式。
	}



	
	
	
\end{enumerate}


	
	\item 
	把传导电子限制在金属内势的一种平均势, 对于下列一维模式（如图)
\[ 
	V(x)=\left\{\begin{array}{cc}
	-V_{0}, & x<0 \\
	0, & x>0
\end{array}\right. 
 \]
	试就:
	\begin{enumerate}
		%\renewcommand{\labelenumi}{\arabic{enumi}.}
		% A(\Alph) a(\alph) I(\Roman) i(\roman) 1(\arabic)
		%设定全局标号series=example	%引用全局变量resume=example
		%[topsep=-0.3em,parsep=-0.3em,itemsep=-0.3em,partopsep=-0.3em]
		%可使用leftmargin调整列表环境左边的空白长度 [leftmargin=0em]
		\item
 $E>0$,
	\item 
	$-V_{0}<E<0$ 
	\end{enumerate}
两种情况计算接近金属表面的传导电子的反射和透射几率。

\banswer{
	设电子从金属内部 $x<0$ 向表面入射。其总能量为 $E$，在两侧的定态波函数取为
	\[
		\psi_{<}=e^{ik_1x}+r e^{-ik_1x},\qquad
		\psi_{>}=t e^{ik_2x},\qquad
		k_1=\frac{\sqrt{2m(E+V_0)}}{\hbar}.
	\]
	
	\begin{enumerate}
		\item 当 $E>0$ 时，$k_2=\sqrt{2mE}/\hbar$。由 $x=0$ 处 $\psi$ 与 $d\psi/dx$ 连续，
		\[
		 1+r=t,\qquad k_1(1-r)=k_2t.
		\]
		解这两个线性方程得
		\[
			r=\frac{k_1-k_2}{k_1+k_2},\qquad
			t=\frac{2k_1}{k_1+k_2}.
		\]
		按概率流之比计算，
		\[
			R=|r|^2=\left(\frac{k_1-k_2}{k_1+k_2}\right)^2,\qquad
			T=\frac{k_2}{k_1}|t|^2
			=\frac{4k_1k_2}{(k_1+k_2)^2},\qquad R+T=1.
		\]

		\item 当 $-V_0<E<0$ 时，令 $\kappa=\sqrt{-2mE}/\hbar$，则真空一侧只有衰减解 $\psi_{>}=t e^{-\kappa x}$。连续条件给出
		\[
			r=\frac{k_1-i\kappa}{k_1+i\kappa},\qquad
			t=\frac{2k_1}{k_1+i\kappa}.
		\]
		因 $|r|=1$ 且衰减波不携带向 $+x$ 的概率流，故
		\[
			R=1,\qquad T=0.
		\]
		电子只能在真空侧穿透约 $1/\kappa$ 的距离，不能形成透射流。
	\end{enumerate}
}



\newpage
\item 
对于一维谐振子, 求消灭算符 $a$ 的本征态, 将其表示成各能量本征态 $|n\rangle$ 的线性叠加。

\banswer{
	设消灭算符的本征值为复数 $\alpha$，本征态为
	\[
		|\alpha\rangle=\sum_{n=0}^{\infty}c_n|n\rangle,\qquad
		a|\alpha\rangle=\alpha|\alpha\rangle.
	\]
	利用 $a|n\rangle=\sqrt n|n-1\rangle$，比较 $|n\rangle$ 的系数可得
	\[
		c_{n+1}\sqrt{n+1}=\alpha c_n,\qquad
		c_n=\frac{\alpha^n}{\sqrt{n!}}c_0.
	\]
	归一化条件给出
	\[
		1=|c_0|^2\sum_{n=0}^{\infty}\frac{|\alpha|^{2n}}{n!}
		=|c_0|^2e^{|\alpha|^2}.
	\]
	略去任意整体相位后 $c_0=e^{-|\alpha|^2/2}$，故消灭算符的归一化本征态（相干态）为
	\[
		\boxed{|\alpha\rangle=e^{-|\alpha|^2/2}
		\sum_{n=0}^{\infty}\frac{\alpha^n}{\sqrt{n!}}|n\rangle}.
	\]
}

\item 
给定 $(\theta, \varphi)$ 方向单位矢量 $\vec{n}=\left(n_{x}, n_{y}, n_{z}\right)=(\sin \theta \cos \varphi, \sin \theta \sin \varphi, \cos \theta)$。求 $\sigma_{n}=\vec{\sigma} \cdot \vec{n}$ 的本征值和本征函数。（取 $\sigma_{z}$ 表象）


\banswer{
	在 $\sigma_z$ 表象中
	\[
		\sigma_n=\vec\sigma\cdot\vec n
		=\begin{pmatrix}
		\cos\theta&e^{-i\varphi}\sin\theta\\
		e^{i\varphi}\sin\theta&-\cos\theta
		\end{pmatrix}.
	\]
	特征方程为
	\[
	 \det(\sigma_n-\lambda I)=(\cos\theta-\lambda)(-\cos\theta-\lambda)
	 -\sin^2\theta=\lambda^2-1=0,
	\]
	所以本征值为 $\lambda_\pm=\pm1$。设本征矢为 $(u,v)^T$，对于 $\lambda=+1$，
	第一行给出 $v/u=e^{i\varphi}(1-\cos\theta)/\sin\theta
	=e^{i\varphi}\tan(\theta/2)$；取 $|u|^2+|v|^2=1$。
	负本征矢与它正交并归一化，得到
	\[
		\chi_+(\theta,\varphi)=
		\begin{pmatrix}\cos(\theta/2)\\e^{i\varphi}\sin(\theta/2)\end{pmatrix},
		\qquad
		\chi_-(\theta,\varphi)=
		\begin{pmatrix}-e^{-i\varphi}\sin(\theta/2)\\\cos(\theta/2)\end{pmatrix}.
	\]
	它们分别满足 $\sigma_n\chi_\pm=\pm\chi_\pm$；相应的自旋投影本征值为 $\pm\hbar/2$。
}

\item 
有一个定域电子（不计及其轨道运动）受到均匀磁场作用，磁场 $B$ 指向正 $x$ 方向, 磁作用势为:
\[
\hat{H}=\frac{e B}{\mu c} \cdot \hat{S}_{x}=\frac{e \hbar B}{2 \mu c} \hat{\sigma}_{x}
\]
设 $t=0$ 时电子的自旋 “向上”, 即 $S_{x}=\hbar / 2$, 求 $t>0$ 时 $\vec{S}$ 的平均值。

\banswer{
	令 $\Omega=eB/(\mu c)$，则
	\[
		\hat H=\frac{\hbar\Omega}{2}\sigma_x.
	\]
	题设初态是 $S_x$ 的正本征态 $|+x\rangle$，也是 $\hat H$ 的本征态。因此
	\[
		|\psi(t)\rangle=e^{-i\Omega t/2}|+x\rangle,
	\]
	时间演化只产生整体相位。于是
	\[
		\boxed{\langle\vec S\rangle_t
		=\left(\frac{\hbar}{2},0,0\right)}.
	\]
}


\end{enumerate}
